\documentclass[a4paper,11pt]{article}
\usepackage[utf8]{inputenc}
\usepackage[T1]{fontenc}
\usepackage[french]{babel}
\usepackage{amsmath, amssymb, amsthm}
\usepackage{graphicx}
\usepackage{tcolorbox}
\usepackage{tikz}
\usetikzlibrary{shapes,arrows,positioning,calc}
\usepackage{listings}
\usepackage{xcolor}
\usepackage{geometry}
\usepackage{hyperref}

\geometry{hmargin=2.5cm,vmargin=2.5cm}

% Configuration des boîtes
\tcbuselibrary{skins,breakable}

\newtcolorbox{conceptbox}[1][]{
  colback=blue!5!white,
  colframe=blue!75!black,
  title={\textbf{Concept Clé}},
  fonttitle=\bfseries,
  enhanced,
  drop shadow,
  breakable,
  #1
}

\newtcolorbox{mathbox}[1][]{
  colback=green!5!white,
  colframe=green!60!black,
  title={\textbf{Fondement Mathématique}},
  fonttitle=\bfseries,
  enhanced,
  drop shadow,
  breakable,
  #1
}

\newtcolorbox{codebox}[1][]{
  colback=gray!5!white,
  colframe=gray!75!black,
  title={\textbf{Implémentation Python}},
  fonttitle=\bfseries,
  breakable,
  #1
}

\newtcolorbox{warningbox}[1][]{
  colback=red!5!white,
  colframe=red!75!black,
  title={\textbf{Attention}},
  fonttitle=\bfseries,
  enhanced,
  #1
}

\newtcolorbox{examplebox}[1][]{
  colback=orange!5!white,
  colframe=orange!75!black,
  title={\textbf{Exemple}},
  fonttitle=\bfseries,
  enhanced,
  breakable,
  #1
}

% Configuration du code
\definecolor{codegreen}{rgb}{0,0.6,0}
\definecolor{codegray}{rgb}{0.5,0.5,0.5}
\definecolor{codepurple}{rgb}{0.58,0,0.82}
\definecolor{backcolour}{rgb}{0.95,0.95,0.92}

\lstdefinestyle{pythonstyle}{
    language=Python,
    backgroundcolor=\color{backcolour},   
    commentstyle=\color{codegreen},
    keywordstyle=\color{magenta},
    numberstyle=\tiny\color{codegray},
    stringstyle=\color{codepurple},
    basicstyle=\ttfamily\footnotesize,
    breakatwhitespace=false,         
    breaklines=true,                 
    captionpos=b,                    
    keepspaces=true,                 
    numbers=left,                    
    numbersep=5pt,                  
    showspaces=false,                
    showstringspaces=false,
    showtabs=false,                  
    tabsize=2
}
\lstset{style=pythonstyle}

\title{\textbf{One-Class SVM} \\ \large Support Vector Machines pour la Détection d'Anomalies}
\author{Cours de Machine Learning Non Supervisé}
\date{\today}

\begin{document}

\maketitle
\tableofcontents
\newpage

\section{Introduction}

Le \textbf{One-Class SVM} (OC-SVM) est une adaptation des Support Vector Machines pour la détection d'anomalies non supervisée.

\begin{conceptbox}
\textbf{Intuition Fondamentale}

Contrairement au SVM classique qui sépare deux classes, le One-Class SVM cherche à :
\begin{itemize}
    \item Entourer les données normales dans un espace transformé
    \item Séparer ces données de l'origine avec une marge maximale
    \item Identifier comme anomalies les points qui tombent en dehors de cette région
\end{itemize}
\end{conceptbox}

\section{Principe de Fonctionnement}

\subsection{Transformation par Noyau}

Comme les SVM classiques, on utilise le \textit{kernel trick} pour projeter les données dans un espace de dimension supérieure où elles deviennent séparables.

\subsection{Séparation de l'Origine}

L'algorithme cherche l'hyperplan qui :
\begin{enumerate}
    \item Sépare la majorité des points de l'origine
    \item Maximise la distance (marge) à l'origine
    \item Tolère quelques violations (points du mauvais côté)
\end{enumerate}

\begin{center}
\begin{tikzpicture}[scale=0.9]
    % Axes
    \draw[->] (-1,0) -- (6,0) node[right] {$x_1$};
    \draw[->] (0,-1) -- (0,5) node[above] {$x_2$};
    
    % Origine
    \node at (0,0) [circle,fill,inner sep=1.5pt,label=below left:Origine] {};
    
    % Points normaux
    \foreach \x/\y in {2/3, 2.5/2.5, 3/3.5, 3.5/2, 4/3, 2.5/3.5, 3/4, 3.5/3}
        \node at (\x,\y) [circle,fill=blue!60,inner sep=2pt] {};
        
    % Anomalie
    \node at (5,1) [circle,fill=red!60,inner sep=2pt,label=right:Anomalie] {};
    
    % Hyperplan
    \draw[thick, dashed] (1,4.5) -- (5,0.5) node[midway,above,sloped] {Frontière};
    
    % Marge
    \draw[->, thick] (1.5,1) -- (2.2,2.2) node[midway, below right] {$\rho$};
    
    % Zone normale
    \fill[blue, opacity=0.1] (1,4.5) -- (5,0.5) -- (6,5) -- (1,5) -- cycle;
    \node at (4,4) [blue] {Zone Normale};
\end{tikzpicture}
\end{center}

\section{Fondements Mathématiques}

\begin{mathbox}
\textbf{Formulation du Problème}

On cherche à minimiser :
$$ \min_{w, \xi, \rho} \frac{1}{2} ||w||^2 + \frac{1}{\nu n} \sum_{i=1}^{n} \xi_i - \rho $$

Sous les contraintes :
$$ (w \cdot \phi(x_i)) \geq \rho - \xi_i, \quad \xi_i \geq 0 $$

Où :
\begin{itemize}
    \item $w$ : vecteur normal à l'hyperplan
    \item $\rho$ : distance de l'hyperplan à l'origine (marge)
    \item $\xi_i$ : variables de relâchement (slack variables)
    \item $\nu \in (0, 1]$ : hyperparamètre de contrôle
    \item $\phi(x)$ : transformation par noyau
\end{itemize}
\end{mathbox}

\subsection{Fonction de Décision}

Pour un nouveau point $x$, la décision est :
$$ f(x) = \text{sign}((w \cdot \phi(x)) - \rho) $$

\begin{itemize}
    \item $f(x) = +1$ : Point normal (dans la région)
    \item $f(x) = -1$ : Anomalie (hors de la région)
\end{itemize}

\section{Hyperparamètres}

\subsection{Le paramètre nu ($\nu$)}

C'est le paramètre le plus important :
\begin{itemize}
    \item \textbf{Borne supérieure} sur la fraction d'anomalies tolérées
    \item \textbf{Borne inférieure} sur la fraction de vecteurs de support
    \item Valeur typique : 0.01 à 0.1 (1\% à 10\%)
\end{itemize}

\begin{examplebox}
Si vous fixez $\nu = 0.05$, vous indiquez que :
\begin{itemize}
    \item Au maximum 5\% des points d'entraînement peuvent être des anomalies
    \item Au minimum 5\% des points seront des vecteurs de support
\end{itemize}
\end{examplebox}

\subsection{Le Noyau (Kernel)}

Le noyau RBF (Radial Basis Function) est le plus utilisé :
$$ K(x, y) = \exp\left(-\gamma ||x - y||^2\right) $$

\subsection{Le paramètre gamma ($\gamma$)}

Contrôle la complexité de la frontière :
\begin{itemize}
    \item $\gamma$ \textbf{élevé} : Frontière complexe, risque de surapprentissage
    \item $\gamma$ \textbf{faible} : Frontière lisse, meilleure généralisation
    \item Valeur par défaut : $\gamma = \frac{1}{n_{features} \times \text{var}(X)}$
\end{itemize}

\section{Avantages et Limites}

\subsection{Avantages}
\begin{itemize}
    \item Peut modéliser des frontières non linéaires complexes
    \item Base théorique solide (théorie des SVM)
    \item Efficace en haute dimension
\end{itemize}

\subsection{Limites}
\begin{itemize}
    \item \textbf{Complexité} : $O(n^2)$ à $O(n^3)$ (lent pour grands datasets)
    \item \textbf{Sensibilité} : Très sensible au choix de $\nu$ et $\gamma$
    \item \textbf{Scalabilité} : Difficile au-delà de 100 000 points
\end{itemize}

\begin{warningbox}
Pour le dataset AI4I (10 000 lignes), One-Class SVM est utilisable mais peut prendre plusieurs secondes à minutes selon les paramètres.
\end{warningbox}

\section{Exemple de Code Python}

\begin{codebox}
\begin{lstlisting}
from sklearn.svm import OneClassSVM
from sklearn.preprocessing import StandardScaler
from sklearn.metrics import classification_report, confusion_matrix
import numpy as np

# 1. PRÉPARATION DES DONNÉES
# La normalisation est CRUCIALE pour One-Class SVM
scaler = StandardScaler()
X_scaled = scaler.fit_transform(X)

# 2. CONFIGURATION DU MODÈLE
ocsvm = OneClassSVM(
    kernel='rbf',      # Noyau gaussien
    gamma='auto',      # gamma = 1 / (n_features * X.var())
    nu=0.03            # On tolère 3% d'anomalies
)

# 3. ENTRAÎNEMENT
print("Entraînement One-Class SVM...")
ocsvm.fit(X_scaled)

# 4. PRÉDICTION
# Retourne +1 (normal) ou -1 (anomalie)
predictions = ocsvm.predict(X_scaled)

# 5. CONVERSION POUR ÉVALUATION
# Nos labels : 0=Normal, 1=Panne
y_pred = np.where(predictions == -1, 1, 0)

# 6. ÉVALUATION
print("Matrice de Confusion :")
print(confusion_matrix(y, y_pred))

print("\nRapport de Classification :")
print(classification_report(y, y_pred, 
                          target_names=['Normal', 'Panne']))

# 7. SCORES BRUTS (pour analyse)
# Plus le score est négatif, plus c'est une anomalie
scores = ocsvm.decision_function(X_scaled)
print(f"\nScore min (anomalie forte) : {scores.min():.4f}")
print(f"Score max (normal fort) : {scores.max():.4f}")
\end{lstlisting}
\end{codebox}

\section{Comparaison avec Autres Méthodes}

\begin{center}
\begin{tabular}{|l|c|c|c|}
\hline
\textbf{Critère} & \textbf{IF} & \textbf{OC-SVM} & \textbf{LOF} \\
\hline
Complexité & $O(n \log n)$ & $O(n^2)$ & $O(n^2)$ \\
Scalabilité & Excellente & Faible & Faible \\
Frontières non-linéaires & Non & Oui & Oui \\
Anomalies locales & Non & Non & Oui \\
Sensibilité paramètres & Faible & Élevée & Moyenne \\
\hline
\end{tabular}
\end{center}

\section{Conclusion}

One-Class SVM est une méthode puissante pour modéliser des frontières de normalité complexes. Elle excelle lorsque :
\begin{itemize}
    \item Les données normales forment une région compacte
    \item La frontière entre normal et anormal est non linéaire
    \item Le dataset n'est pas trop volumineux ($n < 50000$)
\end{itemize}

Pour des datasets plus grands ou des anomalies locales, préférer Isolation Forest ou LOF respectivement.

\end{document}
